\section{距离观测的坐标重构}
\label{sec:smap-distance-reconstruction}

上一节说明，距离最多把点配置确定到合同为止。本节在这个限度内研究重构：完整距离表在什么条件下来自欧氏点、最少需要几维、怎样写出坐标；维数受限时怎样近似；只观测部分距离时形状是否仍然确定；以及测量误差会怎样传到坐标。比较两组坐标时，总要先消去整体平移与正交变换，否则相差一个旋转的两组正确坐标也会被误判为不同。

\subsection{完整距离的欧氏实现}
\label{sec:smap-full-distance}

设$n$个点两两之间的距离全部已知，排成矩阵$\bm D=[d_{i,j}]\in\mathbb R^{n\times n}$。任何来自点的距离表都满足三条基本要求：各项非负，$d_{i,j}=d_{j,i}$，对角元$d_{i,i}=0$。下面始终假定输入满足这三条，并记逐项平方后的矩阵为$\bm\Delta=[\delta_{i,j}]$，$\delta_{i,j}=d_{i,j}^2$。问题是：是否存在$\bm x_1,\ldots,\bm x_n\in\mathbb R^d$，使$\lVert\bm x_i-\bm x_j\rVert_2^2=\delta_{i,j}$对所有$i,j$成立？若存在，称$\bm\Delta$\term{欧氏可实现}（Euclidean realizable），所取的点称为它的一个实现。

\paragraph{从坐标到Gram矩阵}

先从已知坐标出发，看平方距离与内积怎样联系。把各点按行排成$\bm X\in\mathbb R^{n\times d}$，第$i$行为$\bm x_i^\top$，Gram矩阵$\bm G=\bm X\bm X^\top$记录两两内积$g_{i,j}=\bm x_i^\top\bm x_j$。展开平方距离，
\begin{equation}
 \delta_{i,j}=\lVert\bm x_i\rVert_2^2+\lVert\bm x_j\rVert_2^2-2\bm x_i^\top\bm x_j=g_{i,i}+g_{j,j}-2g_{i,j}.
 \label{eq:smap-distance-gram}
\end{equation}
写成矩阵，记$\bm g=(g_{1,1},\ldots,g_{n,n})^\top$、$\bm1$为全一向量，就是$\bm\Delta=\bm g\bm1^\top+\bm1\bm g^\top-2\bm G$。

反方向不能直接解出$\bm G$。距离不依赖原点，内积却依赖：整组点平移后，$\bm\Delta$不变而$\bm G$改变。所以要先约定原点，最自然的选择是质心。令
\begin{equation}
 \bm J=\bm I_n-\frac1n\bm1\bm1^\top ,
 \label{eq:smap-centering-matrix}
\end{equation}
它对称，满足$\bm J^2=\bm J$与$\bm J\bm1=\bm0$，是到$\bm1$的正交补上的正交投影。$\bm J\bm X$从每一行减去各点的平均坐标，所以$\bm J\bm X$的各行就是以质心为原点的坐标。

用$\bm J$左右夹乘$\bm\Delta$。由$\bm J\bm1=\bm0$，$\bm J\bm g\bm1^\top\bm J$与$\bm J\bm1\bm g^\top\bm J$都是零矩阵，只剩
\begin{equation}
 -\frac12\bm J\bm\Delta\bm J=\bm J\bm G\bm J=(\bm J\bm X)(\bm J\bm X)^\top .
 \label{eq:smap-double-centering}
\end{equation}
右端正是中心化坐标的Gram矩阵。这一运算称为\term{双中心化}（double centering）：左乘$\bm J$消去行上的均值项，右乘消去列上的均值项，把平方距离转成以质心为原点的内积。下面记$\bm B=-\frac12\bm J\bm\Delta\bm J$，它对任何对称零对角的$\bm\Delta$都有定义，不必事先知道$\bm\Delta$来自点。

\paragraph{可实现的判据}

若$\bm\Delta$可实现，由式~\eqref{eq:smap-double-centering}，$\bm B$是一个Gram矩阵，因而半正定：$\bm c^\top\bm B\bm c=\lVert(\bm J\bm X)^\top\bm c\rVert_2^2\geq0$。要证明半正定也是充分的，需要确认双中心化没有丢掉信息，即零对角的对称矩阵由它的双中心化唯一确定。

\begin{lemma}[双中心化的核]
\label{lem:smap-double-centering-kernel}
设$\bm H\in\mathbb R^{n\times n}$对称且$\bm J\bm H\bm J=\bm0$，则存在$\bm a\in\mathbb R^n$使$\bm H=\bm a\bm1^\top+\bm1\bm a^\top$。若还有$h_{i,i}=0$对所有$i$成立，则$\bm H=\bm0$。
\end{lemma}
\begin{proof}
记$\bm m=\frac1n\bm H\bm1$，$\mu=\frac1{n^2}\bm1^\top\bm H\bm1$。把$\bm J=\bm I-\frac1n\bm1\bm1^\top$代入$\bm J\bm H\bm J=\bm0$并展开，
\[
 \bm H-\bm1\bm m^\top-\bm m\bm1^\top+\mu\bm1\bm1^\top=\bm0 .
\]
取$\bm a=\bm m-\frac\mu2\bm1$，即得$\bm H=\bm a\bm1^\top+\bm1\bm a^\top$。此时$h_{i,i}=2a_i$，零对角迫使$\bm a=\bm0$，于是$\bm H=\bm0$。
\end{proof}

\begin{theorem}[欧氏可实现的判据]
\label{thm:smap-realizability}
设$\bm\Delta\in\mathbb R^{n\times n}$对称、零对角，$\bm B=-\frac12\bm J\bm\Delta\bm J$。则$\bm\Delta$是某组欧氏点的平方距离矩阵，当且仅当$\bm B$半正定。此时实现所需的最小维数等于$r=\operatorname{rank}\bm B$，并且可以取
\begin{equation}
 \bm Y=\bm Q_r\bm\Lambda_r^{1/2}\in\mathbb R^{n\times r},
 \label{eq:smap-classical-scaling}
\end{equation}
其中$\bm\Lambda_r$是$\bm B$的$r$个正特征值组成的对角矩阵，$\bm Q_r$的列是对应的标准正交特征向量。
\end{theorem}
\begin{proof}
必要性已在上面说明。设$\bm B$半正定，由对称矩阵的谱分解，$\bm B=\bm Q_r\bm\Lambda_r\bm Q_r^\top$，于是式~\eqref{eq:smap-classical-scaling}的$\bm Y$满足$\bm Y\bm Y^\top=\bm B$。由$\bm B\bm1=\bm0$得$\lVert\bm Y^\top\bm1\rVert_2^2=\bm1^\top\bm B\bm1=0$，即$\bm Y$的各行已经中心化，$\bm J\bm Y=\bm Y$。

记$\bm\Delta'$为$\bm Y$各行的平方距离矩阵。由式~\eqref{eq:smap-double-centering}，$-\frac12\bm J\bm\Delta'\bm J=\bm Y\bm Y^\top=\bm B$，所以$\bm H=\bm\Delta-\bm\Delta'$对称、零对角且$\bm J\bm H\bm J=\bm0$。由引理~\ref{lem:smap-double-centering-kernel}，$\bm H=\bm0$，即$\bm Y$实现了$\bm\Delta$。

最后看维数。任何$\mathbb R^k$中的实现，中心化后的Gram矩阵就是$\bm B$，而它是$n\times k$矩阵与其转置之积，秩不超过$k$，故$k\geq r$；上面的构造又在$k=r$时达到。$r=0$时所有点重合。
\end{proof}

这一判据与谱分解构造合称\term{经典多维尺度}（classical multidimensional scaling）。所用的对称矩阵谱分解可参考\BookChapterRef{chap:eigen-spectral-decomposition}{卷二第4章“特征值与谱分解”}。计算步骤是：逐项平方得$\bm\Delta$，双中心化得$\bm B$，求$\bm B$的特征值与特征向量，检查有无负特征值，再按式~\eqref{eq:smap-classical-scaling}写出坐标，最后用重构坐标重新计算距离与原表核对。

\paragraph{坐标的唯一性}

谱分解中特征向量可以变号，重特征值时特征空间内部还可以换基，所以式~\eqref{eq:smap-classical-scaling}给出的坐标并不唯一。下面的定理说明，这些差别恰好是正交变换，并且任何两个同维实现都只差一个合同。

\begin{theorem}[完整距离确定合同类]
\label{thm:smap-full-distance-congruence}
设$\bm x_1,\ldots,\bm x_n$与$\bm x_1',\ldots,\bm x_n'\in\mathbb R^d$的两两距离全部相同，则两组点合同。更具体地，以各自质心为原点后，存在$\bm Q\in\mathrm O(d)$使$\bm x_i'-\overline{\bm x}'=\bm Q(\bm x_i-\overline{\bm x})$对所有$i$成立。
\end{theorem}
\begin{proof}
记$\bm z_i=\bm x_i-\overline{\bm x}$、$\bm z_i'=\bm x_i'-\overline{\bm x}'$。由式~\eqref{eq:smap-double-centering}，两组中心化坐标的Gram矩阵都等于同一个$\bm B$，即$\bm z_i^\top\bm z_j=\bm z_i'^\top\bm z_j'$对所有$i,j$成立。

于是对任意系数$\bm c\in\mathbb R^n$，$\lVert\sum_ic_i\bm z_i\rVert_2^2=\bm c^\top\bm B\bm c=\lVert\sum_ic_i\bm z_i'\rVert_2^2$。取$S=\operatorname{span}\{\bm z_i\}$、$S'=\operatorname{span}\{\bm z_i'\}$，规定$L(\sum_ic_i\bm z_i)=\sum_ic_i\bm z_i'$。若两组系数给出$S$中同一向量，它们的差$\bm c$使$\sum_ic_i\bm z_i=\bm0$，由上式$\sum_ic_i\bm z_i'=\bm0$，所以$L$是良定义的线性映射，并且保持长度。

保长的线性映射是单射，所以$\dim S'=\dim S$，正交补$S^\perp$与$S'^\perp$维数也相同。在两个正交补中各取一组标准正交基，把前者逐个送到后者，就把$L$延拓为$\mathbb R^d$上保长的线性映射$\bm Q$，即正交矩阵，并且$\bm Q\bm z_i=\bm z_i'$。
\end{proof}

这一定理兑现了第~\ref{sec:smap-identifiability}节的论断：对完整距离表，合同是唯一的不确定性。$\bm\Lambda_r$的特征值互不相同时，式~\eqref{eq:smap-classical-scaling}中特征向量的选取只差各列的符号，对应沿各坐标轴的反射；有重特征值时，相应的坐标轴可以在特征空间内任意旋转。

\paragraph{实现与不可实现的例子}

\begin{example}[三个点的重构]
\label{ex:smap-three-point-mds}
取$\bm x_1=(-1,0)^\top$、$\bm x_2=(0,1)^\top$、$\bm x_3=(1,0)^\top$，平方距离为$\delta_{1,2}=\delta_{2,3}=2$、$\delta_{1,3}=4$。记一般的$\delta_{1,2}=\delta_{2,3}=a$、$\delta_{1,3}=b$，由对称性可以直接读出$\bm B$的特征向量：$\bm1$对应特征值$0$；对$\bm v=(1,0,-1)^\top/\sqrt2$，由$\bm J\bm v=\bm v$得$\bm v^\top\bm B\bm v=-\frac12\bm v^\top\bm\Delta\bm v=b/2$；对$\bm w=(1,-2,1)^\top/\sqrt6$，同样算得$\bm w^\top\bm B\bm w=(4a-b)/6$。

代入$a=2$、$b=4$，正特征值为$2$与$2/3$，所以精确重构恰好需要二维。由式~\eqref{eq:smap-classical-scaling}，重构坐标的两列为$\sqrt2\,\bm v=(1,0,-1)^\top$与$\sqrt{2/3}\,\bm w=(1,-2,1)^\top/3$，三个点是$(1,1/3)$、$(0,-2/3)$、$(-1,1/3)$。它们两两的平方距离为$1+1=2$、$1+1=2$、$4+0=4$，与原表一致。原配置平移到质心后为$(-1,-1/3)$、$(0,2/3)$、$(1,-1/3)$，重构结果恰是它旋转$180^\circ$的结果，与原配置合同；特征向量取相反的符号，就得到中心化的原配置本身。
\end{example}

\begin{example}[违反三角不等式的距离表]
\label{ex:smap-nonrealizable}
取$d_{1,2}=d_{2,3}=1$、$d_{1,3}=3$，即$a=1$、$b=9$。按上例的公式，$\bm B$的特征值为$9/2$、$(4-9)/6=-5/6$与$0$，出现负特征值，所以这张表在任何维数的欧氏空间中都不可实现。这与直接观察一致：$3>1+1$违反三角不等式。负特征值的特征向量$(1,-2,1)^\top$指出问题出在中间点相对两端的位置上。
\end{example}

\subsection{低维坐标重构}
\label{sec:smap-low-dimension}

实际中常常指定维数$k$，例如要在平面上画出点的分布，而$\bm B$的秩可能大于$k$；含误差的距离表还可能让$\bm B$出现负特征值，根本不可实现。这时只能求近似坐标$\bm Y\in\mathbb R^{n\times k}$。近似必须先说明比较什么，本小节先以Gram矩阵为比较对象，再看距离上的误差。

\paragraph{Gram矩阵的最优截断}

以Frobenius范数比较Gram矩阵，问题是在秩不超过$k$的半正定矩阵中找离$\bm B$最近的一个。$\bm B$半正定时，它的谱分解就是奇异值分解，截断到前$k$个特征值即得最优解，这是最佳低秩逼近定理的直接推论，可参考\BookChapterRef{chap:matrix-approximation}{卷二第15章“矩阵逼近”}。$\bm B$有负特征值时还要多一步：负谱部分对任何半正定近似都只会增加误差。

\begin{theorem}[半正定低秩近似]
\label{thm:smap-psd-truncation}
设$\bm B$对称，特征值$\lambda_1\geq\cdots\geq\lambda_n$，特征向量为标准正交的$\bm q_1,\ldots,\bm q_n$。令$p$为正特征值的个数，$k'=\min(k,p)$，$\bm B_k=\sum_{l\leq k'}\lambda_l\bm q_l\bm q_l^\top$。则对一切秩不超过$k$的半正定矩阵$\bm X$，
\[
 \lVert\bm B-\bm X\rVert_{\mathrm F}^2\geq\lVert\bm B-\bm B_k\rVert_{\mathrm F}^2=\sum_{l=k'+1}^{p}\lambda_l^2+\sum_{\lambda_l<0}\lambda_l^2 .
\]
\end{theorem}
\begin{proof}
把$\bm B$分成正部与负部：
\begin{equation*}
\bm B_+=\sum_{\lambda_l>0}\lambda_l\bm q_l\bm q_l^\top,\qquad\bm B_-=\sum_{\lambda_l<0}(-\lambda_l)\bm q_l\bm q_l^\top,
\end{equation*}
二者都半正定，$\bm B=\bm B_+-\bm B_-$，且由特征向量正交，$\bm B_+\bm B_-=\bm0$。对半正定的$\bm X$展开，
\[
 \lVert\bm B_+-\bm X-\bm B_-\rVert_{\mathrm F}^2
 =\lVert\bm B_+-\bm X\rVert_{\mathrm F}^2+\lVert\bm B_-\rVert_{\mathrm F}^2+2\operatorname{tr}(\bm X\bm B_-),
\]
这里用了$\operatorname{tr}(\bm B_+\bm B_-)=0$。两个半正定矩阵之积的迹非负：写$\bm B_-=\bm C\bm C^\top$，则$\operatorname{tr}(\bm X\bm B_-)=\operatorname{tr}(\bm C^\top\bm X\bm C)\geq0$。

所以$\lVert\bm B-\bm X\rVert_{\mathrm F}^2\geq\lVert\bm B_+-\bm X\rVert_{\mathrm F}^2+\lVert\bm B_-\rVert_{\mathrm F}^2$。$\bm B_+$半正定，由最佳低秩逼近定理，秩不超过$k$时$\lVert\bm B_+-\bm X\rVert_{\mathrm F}^2\geq\sum_{l=k'+1}^p\lambda_l^2$，等号在$\bm X=\bm B_k$时取到；而$\bm X=\bm B_k$时$\operatorname{tr}(\bm B_k\bm B_-)=0$，上面的不等式都取等号。
\end{proof}

于是低维坐标取$\bm Y_k=\bm Q_{k'}\bm\Lambda_{k'}^{1/2}$，即只保留前$k'$个正特征值对应的列。它仍然中心化，因为每个保留的特征向量都与$\bm1$正交。负特征值被舍弃，舍弃量$\sum_{\lambda_l<0}\lambda_l^2$应与测量误差的大小比较：若它远大于误差能解释的范围，问题出在数据或模型上，而不是维数选小了。

\paragraph{距离上的误差}

Gram矩阵的误差并不直接说明距离差了多少，需要单独计算。设$\bm\Delta$可实现，$\bm Y=\bm Q_r\bm\Lambda_r^{1/2}$是它的中心化实现。$\bm Y_k$恰好是$\bm Y$的前$k$列，相当于把每个点正交投影到前$k$个坐标轴张成的子空间上。正交投影不增加长度，所以每一对点的距离只会缩短：$\lVert\widetilde{\bm y}_i-\widetilde{\bm y}_j\rVert_2\leq d_{i,j}$。

缩短的总量可以精确算出。对中心化的点，$\sum_{i<j}\lVert\bm y_i-\bm y_j\rVert_2^2=n\sum_i\lVert\bm y_i\rVert_2^2=n\operatorname{tr}\bm B$，这只需展开平方并用$\sum_i\bm y_i=\bm0$。对截断坐标同样有$\sum_{i<j}\widetilde\delta_{i,j}=n\operatorname{tr}\bm B_k$，相减得
\[
 \sum_{i<j}\bigl(\delta_{i,j}-\widetilde\delta_{i,j}\bigr)=n\sum_{l>k}\lambda_l .
\]
注意这里是特征值之和，而Gram误差的平方是特征值的平方和。两种误差都由舍弃的谱决定，但不是同一个量。

\begin{example}[三个点的一维近似]
\label{ex:smap-one-dimensional}
回到例~\ref{ex:smap-three-point-mds}，取$k=1$，只保留特征值$2$。一维坐标为$(1,0,-1)$，两端距离仍为$2$，中间点到两端的距离由$\sqrt2$缩成$1$。Gram误差为舍弃的特征值$2/3$；平方距离的总缩短量为$(2-1)+(2-1)+0=2$，恰为$n\lambda_2=3\cdot2/3$。
\end{example}

\paragraph{最优性依赖目标}

谱截断在Gram矩阵的Frobenius误差下最优，不等于在一切几何准则下最优。常用的另一种准则直接比较距离，称为应力：$\sum_{i<j}(\lVert\widetilde{\bm y}_i-\widetilde{\bm y}_j\rVert_2-d_{i,j})^2$。在上例中，截断坐标的应力为$2(\sqrt2-1)^2=2(3-2\sqrt2)\approx0.343$。

改取一维坐标$(-s,0,s)$，应力为$2(s-\sqrt2)^2+(2s-2)^2$，对$s$求导令其为零，得$s=(2+\sqrt2)/3\approx1.138$。代回去，$s-\sqrt2=(2-2\sqrt2)/3$，$2s-2=(2\sqrt2-2)/3$，应力为$3\cdot4(\sqrt2-1)^2/9=4(3-2\sqrt2)/3\approx0.229$，比截断坐标小。截断坐标把两端拉得太远、中间缩得太多，而应力准则愿意把各对距离的误差摊平。应力关于坐标不是凸的，一般需要迭代优化，而不能由一次谱分解得到。

\subsection{部分距离的重构}
\label{sec:smap-partial-distance}

许多测量只能给出部分点对的距离，例如传感器只能测到通信范围内的邻居。把已观测的点对看作边，就得到\term{观测图}（observation graph）$\mathcal G=(V,E)$；图的基本概念可参考\BookChapterRef{chap:graph-objects}{卷三第1章“图与连接结构”}。观测图不连通时，不同分量之间没有任何约束，各分量可以独立地移动，这是显然的；下面两个例子说明，连通也远远不够。

\paragraph{连通而形状不定}

在平面上取三个点，只观测$d_{1,2}=d_{2,3}=1$。观测图是一条路径，是连通的。但$\bm x_2$处的夹角可以任意选取：$\bm x_1$与$\bm x_3$都在以$\bm x_2$为圆心的单位圆上，它们的距离可以是$[0,2]$中的任何值。不同的夹角给出互不合同的配置，所以连通不保证形状确定。

\paragraph{局部确定而全局不唯一}

为了区分两种“确定”，先给出定义。设$\bm x_1,\ldots,\bm x_n\in\mathbb R^d$满足观测图$\mathcal G$上的全部距离，称它是一个实现。

\begin{definition}[局部刚性与全局刚性]
\label{def:smap-rigidity}
若存在$\eta>0$，使得每个满足同一组观测距离、且$\lVert\bm x_i'-\bm x_i\rVert_2<\eta$对所有$i$成立的实现都与原实现合同，则称原实现\term{局部刚性}（locally rigid）。若满足观测距离的每个实现都与它合同，则称它\term{全局刚性}（globally rigid）。
\end{definition}

取平面上$\bm x_1=(0,0)^\top$、$\bm x_2=(2,0)^\top$、$\bm x_3=(1,1)^\top$、$\bm x_4=(1,-2)^\top$，观测除$\{3,4\}$外的五对距离：$d_{1,2}=2$，$d_{1,3}=d_{2,3}=\sqrt2$，$d_{1,4}=d_{2,4}=\sqrt5$。这是两个共用边$\{1,2\}$的三角形。

先看局部刚性。三角形的三边确定它到合同为止，所以$\bm x_1,\bm x_2,\bm x_3$整体只能作合同变换；在此基础上，$\bm x_4$到$\bm x_1,\bm x_2$的距离给定，它只能落在两个圆的两个交点之一，即直线$\bm x_1\bm x_2$的两侧各一个。两个交点相距$4$，小范围的扰动不能从一个跳到另一个，所以原实现局部刚性。

再看全局。把$\bm x_4$关于直线$\bm x_1\bm x_2$反射到$(1,2)^\top$，五个观测距离都不变，未观测的$d_{3,4}$却由$3$变成$1$。两个实现互不合同，所以原实现不是全局刚性的。局部确定只排除了附近的连续形变，不排除远处另有一个离散的解。

\paragraph{缺边、不相容与含噪}

固定维数$d$，并按第~\ref{sec:smap-identifiability}节的方式约定参照以消去合同自由度之后，部分距离问题有三种不同的困难。缺边可能使解不唯一，如上面的路径和两个三角形；观测不相容时根本无解，例如三个点的三条观测边长为$1$、$1$、$3$；观测含噪时一般也无精确解，只能在已观测的边上度量偏差。

含噪时自然的做法是只对已观测的边建立残差，例如
\[
 \sum_{\{i,j\}\in E}\bigl(\lVert\bm x_i-\bm x_j\rVert_2-d_{i,j}\bigr)^2 ,
\]
或用平方距离之差代替距离之差。未观测的点对不进入目标，因为对它们没有任何数据。这类目标关于坐标不是凸的，可能有多个局部极小，对应的正是上面的离散多解；其求解属于迭代优化。

\paragraph{不能把缺失填成零}

把缺失的距离填零，再套用第~\ref{sec:smap-full-distance}节的谱重构，是常见的错误。零距离不是“未知”，而是断言两点重合。在路径的例子中填$d_{1,3}=0$，得到完整表$\delta_{1,2}=\delta_{2,3}=1$、$\delta_{1,3}=0$，按例~\ref{ex:smap-three-point-mds}的公式，$a=1$、$b=0$，$\bm B$的特征值为$(4a-b)/6=2/3$与$b/2=0$，表是可实现的：$\bm x_1=\bm x_3$，链被完全折叠。

谱重构会毫无警告地给出这个配置，好像形状已经确定，实际上它只是无穷多个可行形状中被错误的数据选中的一个。填入其他猜测值同样会选出某个特定形状。所以缺失距离只能作为“无约束”处理，或者在明确的附加假设下另行估计，而不能伪装成观测。

\subsection{坐标重构的稳定性}
\label{sec:smap-reconstruction-stability}

距离表总带有测量误差。本小节先研究完整表上的误差怎样经过平方、双中心化和谱分解传到坐标，再说明部分距离为什么需要另作分析。所用的特征值与特征子空间扰动界取自\BookChapterRef{chap:condition-and-spectral-perturbation}{卷二第7章“条件数与谱扰动”}：对称矩阵受对称扰动$\bm E$时，每个按降序配对的特征值至多移动$\lVert\bm E\rVert_2$（Weyl界）；若前$k$个特征值与其余特征值之间有间隔$g$，且$\lVert\bm E\rVert_2\leq g/2$，前$k$个特征向量张成的子空间最大主夹角的正弦不超过$2\lVert\bm E\rVert_2/g$（Davis--Kahan界）。

\paragraph{从距离误差到Gram误差}

设观测的平方距离表为$\widetilde{\bm\Delta}=\bm\Delta+\bm E$，$\bm E$对称、零对角。双中心化是线性的，$\widetilde{\bm B}-\bm B=-\frac12\bm J\bm E\bm J$。$\bm J$是正交投影，谱范数为$1$，所以
\[
 \lVert\widetilde{\bm B}-\bm B\rVert_2\leq\frac12\lVert\bm E\rVert_2 .
\]
注意误差是加在平方距离上的：若距离本身有误差$\varepsilon_{i,j}$，平方距离的误差为$2d_{i,j}\varepsilon_{i,j}+\varepsilon_{i,j}^2$，长距离上的同样误差会被放大。

\begin{theorem}[完整距离表的谱扰动]
\label{thm:smap-mds-perturbation}
在上述记号下，设$\lambda_1\geq\cdots\geq\lambda_n$与$\widetilde\lambda_1\geq\cdots\geq\widetilde\lambda_n$分别是$\bm B$与$\widetilde{\bm B}$的特征值，则$|\widetilde\lambda_l-\lambda_l|\leq\frac12\lVert\bm E\rVert_2$。若$g=\lambda_k-\lambda_{k+1}>0$且$\lVert\bm E\rVert_2\leq g$，则两者前$k$个特征向量张成的子空间满足$\sin\theta_{\max}\leq\lVert\bm E\rVert_2/g$。
\end{theorem}
\begin{proof}
对$\bm B$与扰动$\widetilde{\bm B}-\bm B$分别应用Weyl界与Davis--Kahan界，扰动的谱范数不超过$\frac12\lVert\bm E\rVert_2$；条件$\lVert\bm E\rVert_2\leq g$保证$\frac12\lVert\bm E\rVert_2\leq g/2$，代入$2\cdot\frac12\lVert\bm E\rVert_2/g$即得。
\end{proof}

\paragraph{谱隙与小特征值}

定理分两层说明误差的去向。特征值本身稳定，误差至多为$\frac12\lVert\bm E\rVert_2$。坐标轴的方向取决于谱隙：相邻特征值接近时，很小的误差就能让坐标轴在相应的特征空间内大幅转动，这与重特征值时坐标轴本来就不唯一是同一现象的两面。

坐标的长度还要再经过开方。第$l$列坐标按$\sqrt{\lambda_l}$缩放，而
\[
 \bigl|\sqrt{\widetilde\lambda_l}-\sqrt{\lambda_l}\bigr|=\frac{|\widetilde\lambda_l-\lambda_l|}{\sqrt{\widetilde\lambda_l}+\sqrt{\lambda_l}}\leq\frac{|\widetilde\lambda_l-\lambda_l|}{\sqrt{\lambda_l}} ,
\]
特征值越小，相对误差越大。更严重的是，误差与小特征值同量级时，$\widetilde\lambda_l$可能变成负数，这个方向就整个丢失了。

\begin{example}[几乎共线的三个点]
\label{ex:smap-near-collinear}
取$\bm x_1=(-1,0)^\top$、$\bm x_2=(0,h)^\top$、$\bm x_3=(1,0)^\top$，$h=0.1$。这时$a=\delta_{1,2}=\delta_{2,3}=1.01$，$b=\delta_{1,3}=4$，由例~\ref{ex:smap-three-point-mds}的公式，特征值为$b/2=2$与$(4a-b)/6=2h^2/3\approx0.0067$。

设$d_{1,3}$被测成$2.02$，只有$1\%$的误差，于是$\widetilde\delta_{1,3}=4.0804$。第二个特征值变为$(4.04-4.0804)/6\approx-0.0067$，由正变负。事实上$2.02>2\sqrt{1.01}\approx2.0100$已违反三角不等式，表在任何维数中都不可实现。第一个特征值只从$2$变为$2.0402$，相对变化约$2\%$；纵向那个高为$0.1$的坐标则完全丢失。这里$\bm E$只有$(1,3)$、$(3,1)$两项$0.0804$，$\lVert\bm E\rVert_2=0.0804$，两个特征值的变化都在Weyl界$0.0402$之内。
\end{example}

\paragraph{比较坐标前消去合同}

评价重构误差时，不能直接比较$\widetilde{\bm Y}$与真实坐标$\bm X$的各个分量：即使距离完全精确，重构也只确定到合同为止。正确的做法是先把两组坐标都平移到质心，再求
\[
 \min_{\bm Q\in\mathrm O(d)}\lVert\widetilde{\bm Y}-\bm X\bm Q\rVert_{\mathrm F},
\]
它的解由第~\ref{sec:smap-rotation}节的正交对齐给出。这里必须允许反射，因为距离同样无法辨认反射。

\paragraph{部分距离的稳定性}

只观测部分距离时，没有完整的$\bm B$可以分析，稳定性取决于观测约束本身是否足够，不能由某张完整表的谱隙推出。仍以两个共边三角形为例，把$\bm x_4$改为$(1,-\eta)^\top$，$\eta>0$很小，第四个点贴近公共边。由$d_{1,4}^2=1+\eta^2$，若$d_{1,4}^2$与$d_{2,4}^2$都有误差$\varepsilon$，$\eta^2$就改变$\varepsilon$，$\eta$约改变$\varepsilon/(2\eta)$，放大倍数随$\eta\to0$无界增长。

同时，$\bm x_4$在公共边两侧的两个解只相距$2\eta$，误差稍大就无法判断它在哪一侧。这些现象只与观测图在这组点上的约束几何有关，而与真实配置的完整距离表是否有良好的谱隙无关：真实配置的完整表在这个例子中完全正常。
